% ═══════════════════════════════════════════════════════════════
% MUSTERLÖSUNG: Oxide – Formeln und Reaktionsgleichungen
% Chemie Klasse 8 | Lehrerversion mit schrittweisen Lösungen
% ═══════════════════════════════════════════════════════════════

\documentclass[11pt, a4paper]{article}

\usepackage[utf8]{inputenc}
\usepackage[T1]{fontenc}
\usepackage[ngerman]{babel}

\usepackage{helvet}
\renewcommand{\familydefault}{\sfdefault}

\usepackage{microtype}
\usepackage[margin=2cm, top=2cm, bottom=1.5cm]{geometry}
\usepackage{enumitem}
\usepackage{tabularx}
\usepackage{booktabs}
\usepackage{array}

\usepackage{fancyhdr}
\usepackage{lastpage}

\usepackage{xcolor}
\usepackage{tcolorbox}
\tcbuselibrary{skins}

% ═══════════════════════════════════════════════════════════════
% FARBEN
% ═══════════════════════════════════════════════════════════════

\definecolor{chemie}{HTML}{2a7a4b}
\definecolor{hellgrau}{HTML}{F5F5F5}
\definecolor{mittelgrau}{HTML}{888888}
\definecolor{dunkelgrau}{HTML}{444444}
\definecolor{loesung}{HTML}{C62828}

\colorlet{fachfarbe}{chemie}

% ═══════════════════════════════════════════════════════════════
% KOPF- UND FUSSZEILE
% ═══════════════════════════════════════════════════════════════

\pagestyle{fancy}
\fancyhf{}
\renewcommand{\headrulewidth}{0.4pt}
\renewcommand{\footrulewidth}{0pt}
\lhead{\small Chemie Klasse 8 | Reaktionsgleichungen}
\rhead{\small\textcolor{loesung}{\textbf{MUSTERLÖSUNG}}}
\cfoot{\small\textcolor{mittelgrau}{Seite \thepage{} von \pageref{LastPage}}}

% ═══════════════════════════════════════════════════════════════
% BOXEN
% ═══════════════════════════════════════════════════════════════

\newtcolorbox{infobox}[1][]{
    colback=hellgrau,
    colframe=hellgrau,
    boxrule=0pt,
    arc=0pt,
    left=8pt, right=8pt, top=6pt, bottom=6pt,
    borderline west={3pt}{0pt}{fachfarbe},
    #1
}

\newtcolorbox{loesungbox}{
    colback=white,
    colframe=loesung,
    boxrule=1pt,
    arc=0pt,
    left=8pt, right=8pt, top=6pt, bottom=6pt
}

\newtcolorbox{schrittbox}{
    colback=hellgrau,
    colframe=hellgrau,
    boxrule=0pt,
    arc=0pt,
    left=8pt, right=8pt, top=6pt, bottom=6pt,
    borderline west={3pt}{0pt}{mittelgrau}
}

% ═══════════════════════════════════════════════════════════════
% BEFEHLE
% ═══════════════════════════════════════════════════════════════

\newcommand{\lsg}[1]{\textcolor{loesung}{\textbf{#1}}}

\newcommand{\abschnitt}[1]{%
    \vspace{10pt}%
    \noindent{\large\bfseries\textcolor{fachfarbe}{#1}}%
    \vspace{6pt}%
    \nopagebreak[4]%
}

\setlength{\parindent}{0pt}
\setlength{\parskip}{4pt}

\newcolumntype{L}[1]{>{\raggedright\arraybackslash}p{#1}}
\newcolumntype{C}[1]{>{\centering\arraybackslash}p{#1}}

% ═══════════════════════════════════════════════════════════════
% DOKUMENT
% ═══════════════════════════════════════════════════════════════

\begin{document}

% ═══════════════════════════════════════════════════════════════
% HEADER
% ═══════════════════════════════════════════════════════════════

\begin{center}
    {\Large\bfseries\textcolor{fachfarbe}{Übung: Oxide}}\\[0.3em]
    {\normalsize\textcolor{loesung}{\textbf{Musterlösung mit schrittweisen Erklärungen}}}
\end{center}
\vspace{-0.3em}
\hrule
\vspace{0.8em}

% ═══════════════════════════════════════════════════════════════
% AUFGABE A - LÖSUNGEN
% ═══════════════════════════════════════════════════════════════

\abschnitt{Aufgabe A: Oxidformeln aufstellen}

\begin{center}
\renewcommand{\arraystretch}{1.6}
\begin{tabular}{|C{0.5cm}|L{3.5cm}|C{2cm}|C{2.5cm}|L{5cm}|}
\hline
\textbf{\#} & \textbf{Name} & \textbf{Wertigkeit} & \textbf{Formel} & \textbf{Erklärung} \\
\hline
1 & Bariumoxid & Ba (II) & \lsg{BaO} & II + II $\rightarrow$ kürzen: 1:1 \\
\hline
2 & Natriumoxid & Na (I) & \lsg{Na$_2$O} & I + II $\rightarrow$ kreuzen: 2:1 \\
\hline
3 & Chromoxid & Cr (III) & \lsg{Cr$_2$O$_3$} & III + II $\rightarrow$ kreuzen: 2:3 \\
\hline
4 & Titandioxid & Ti (IV) & \lsg{TiO$_2$} & IV + II $\rightarrow$ kürzen: 1:2 \\
\hline
5 & Manganoxid & Mn (IV) & \lsg{MnO$_2$} & IV + II $\rightarrow$ kürzen: 1:2 \\
\hline
6 & Zinnoxid & Sn (IV) & \lsg{SnO$_2$} & IV + II $\rightarrow$ kürzen: 1:2 \\
\hline
7 & Silberoxid & Ag (I) & \lsg{Ag$_2$O} & I + II $\rightarrow$ kreuzen: 2:1 \\
\hline
\end{tabular}
\end{center}

\vspace{0.8em}

% ═══════════════════════════════════════════════════════════════
% AUFGABE B - LÖSUNGEN
% ═══════════════════════════════════════════════════════════════

\abschnitt{Aufgabe B: Reaktionsgleichungen (ausgeglichen)}

\begin{center}
\renewcommand{\arraystretch}{1.8}
\begin{tabular}{|C{0.5cm}|L{3cm}|L{10cm}|}
\hline
\textbf{\#} & \textbf{Reaktion} & \textbf{Ausgeglichene Formelgleichung} \\
\hline
1 & Barium & \lsg{2 Ba + O$_2$ $\rightarrow$ 2 BaO} \\
\hline
2 & Natrium & \lsg{4 Na + O$_2$ $\rightarrow$ 2 Na$_2$O} \\
\hline
3 & Chrom & \lsg{4 Cr + 3 O$_2$ $\rightarrow$ 2 Cr$_2$O$_3$} \\
\hline
4 & Titan & \lsg{Ti + O$_2$ $\rightarrow$ TiO$_2$} \\
\hline
5 & Mangan & \lsg{Mn + O$_2$ $\rightarrow$ MnO$_2$} \\
\hline
6 & Zinn & \lsg{Sn + O$_2$ $\rightarrow$ SnO$_2$} \\
\hline
7 & Silber & \lsg{4 Ag + O$_2$ $\rightarrow$ 2 Ag$_2$O} \\
\hline
\end{tabular}
\end{center}

\vspace{1em}

% ═══════════════════════════════════════════════════════════════
% SCHRITTWEISE LÖSUNGEN
% ═══════════════════════════════════════════════════════════════

\abschnitt{Schrittweise Lösungen (für Tafelarbeit)}

\begin{loesungbox}
\textbf{Beispiel 1: Bariumoxid (einfach)}

\textbf{Schritt 1:} Wertigkeiten notieren\\
Ba (II) \quad O (II)

\textbf{Schritt 2:} Kreuzregel anwenden\\
II : II = 1 : 1 (kürzen!) $\rightarrow$ \textbf{BaO}

\textbf{Schritt 3:} Reaktionsgleichung aufstellen\\
Ba + O$_2$ $\rightarrow$ BaO \quad (noch nicht ausgeglichen!)

\textbf{Schritt 4:} Ausgleichen\\
Links: 1 Ba, 2 O \quad Rechts: 1 Ba, 1 O\\
$\rightarrow$ \textbf{2 Ba + O$_2$ $\rightarrow$ 2 BaO}
\end{loesungbox}

\vspace{0.5em}

\begin{loesungbox}
\textbf{Beispiel 2: Natriumoxid (Kreuzregel)}

\textbf{Schritt 1:} Wertigkeiten notieren\\
Na (I) \quad O (II)

\textbf{Schritt 2:} Kreuzregel anwenden\\
I : II $\rightarrow$ kreuzen: Na$_2$O$_1$ = \textbf{Na$_2$O}

\textbf{Schritt 3:} Reaktionsgleichung aufstellen\\
Na + O$_2$ $\rightarrow$ Na$_2$O \quad (noch nicht ausgeglichen!)

\textbf{Schritt 4:} Ausgleichen\\
Links: 1 Na, 2 O \quad Rechts: 2 Na, 1 O\\
$\rightarrow$ Rechts verdoppeln: 2 Na$_2$O (4 Na, 2 O)\\
$\rightarrow$ Links anpassen: 4 Na\\
$\rightarrow$ \textbf{4 Na + O$_2$ $\rightarrow$ 2 Na$_2$O}
\end{loesungbox}

\vspace{0.5em}

\begin{loesungbox}
\textbf{Beispiel 3: Chromoxid (anspruchsvoll)}

\textbf{Schritt 1:} Wertigkeiten notieren\\
Cr (III) \quad O (II)

\textbf{Schritt 2:} Kreuzregel anwenden\\
III : II $\rightarrow$ kreuzen: Cr$_2$O$_3$ = \textbf{Cr$_2$O$_3$}

\textbf{Schritt 3:} Reaktionsgleichung aufstellen\\
Cr + O$_2$ $\rightarrow$ Cr$_2$O$_3$ \quad (noch nicht ausgeglichen!)

\textbf{Schritt 4:} Ausgleichen\\
Rechts: 2 Cr, 3 O\\
$\rightarrow$ Problem: 3 O rechts, aber O$_2$ links (immer gerade!)\\
$\rightarrow$ Rechts verdoppeln: 2 Cr$_2$O$_3$ (4 Cr, 6 O)\\
$\rightarrow$ Links: 4 Cr + 3 O$_2$\\
$\rightarrow$ \textbf{4 Cr + 3 O$_2$ $\rightarrow$ 2 Cr$_2$O$_3$}
\end{loesungbox}

\newpage

\begin{loesungbox}
\textbf{Beispiel 4: Silberoxid (Übergang zur Reduktion)}

\textbf{Schritt 1:} Wertigkeiten notieren\\
Ag (I) \quad O (II)

\textbf{Schritt 2:} Kreuzregel anwenden\\
I : II $\rightarrow$ kreuzen: Ag$_2$O$_1$ = \textbf{Ag$_2$O}

\textbf{Schritt 3:} Reaktionsgleichung aufstellen\\
Ag + O$_2$ $\rightarrow$ Ag$_2$O \quad (noch nicht ausgeglichen!)

\textbf{Schritt 4:} Ausgleichen\\
Rechts: 2 Ag, 1 O\\
$\rightarrow$ Rechts verdoppeln: 2 Ag$_2$O (4 Ag, 2 O)\\
$\rightarrow$ Links: 4 Ag + O$_2$\\
$\rightarrow$ \textbf{4 Ag + O$_2$ $\rightarrow$ 2 Ag$_2$O}

\vspace{0.5em}
\textcolor{fachfarbe}{\textbf{Überleitung:}} Diese Reaktion lässt sich umkehren!\\
Beim Erhitzen von Silberoxid entsteht wieder Silber:\\
\textbf{2 Ag$_2$O $\rightarrow$ 4 Ag + O$_2$} \quad (Reduktion)
\end{loesungbox}

\vspace{1em}

% ═══════════════════════════════════════════════════════════════
% HÄUFIGE FEHLER
% ═══════════════════════════════════════════════════════════════

\abschnitt{Häufige Fehler}

\begin{schrittbox}
\textbf{Fehler 1: Kreuzregel ohne Kürzen}\\
\textit{Falsch:} Ba (II) + O (II) $\rightarrow$ Ba$_2$O$_2$\\
\textit{Richtig:} Kürzen! $\rightarrow$ BaO

\vspace{6pt}

\textbf{Fehler 2: Sauerstoff als O statt O$_2$}\\
\textit{Falsch:} 2 Ba + 2 O $\rightarrow$ 2 BaO\\
\textit{Richtig:} Sauerstoff ist \textbf{immer} O$_2$ (Molekül)!

\vspace{6pt}

\textbf{Fehler 3: Ungerade Sauerstoffzahl}\\
\textit{Problem:} Cr$_2$O$_3$ hat 3 Sauerstoff, aber O$_2$ ist immer gerade\\
\textit{Lösung:} Produkt verdoppeln, dann ausgleichen

\vspace{6pt}

\textbf{Fehler 4: Index statt Vorfaktor ändern}\\
\textit{Falsch:} Ba + O$_2$ $\rightarrow$ BaO$_2$ (Index geändert!)\\
\textit{Richtig:} 2 Ba + O$_2$ $\rightarrow$ 2 BaO (Vorfaktor geändert)
\end{schrittbox}

\vspace{1em}

% ═══════════════════════════════════════════════════════════════
% HINTERGRUND
% ═══════════════════════════════════════════════════════════════

\abschnitt{Hintergrundwissen (optional)}

\begin{infobox}
\small
\begin{tabular}{L{2.5cm}L{11cm}}
\textbf{Bariumoxid} & Feuerwerk (grüne Flammenfarbe), Trocknungsmittel \\
\textbf{Natriumoxid} & Glasherstellung (Soda-Kalk-Glas) \\
\textbf{Chromoxid} & Pigment ``Chromgrün'', Schleifmittel, Katalysator \\
\textbf{Titandioxid} & Weißpigment Nr. 1 weltweit (Farben, Kunststoff, Papier), Sonnencreme (UV-Schutz), Zahnpasta, Lebensmittelfarbe E171 \\
\textbf{Manganoxid} & ``Braunstein'' in Batterien (Zink-Kohle, Alkaline) \\
\textbf{Zinnoxid} & Keramik-Glasuren (weiß, deckend), Glas-Polierung \\
\textbf{Silberoxid} & Angelaufenes Silber (schwarze Schicht), Knopfzellen \\
\end{tabular}
\end{infobox}

\vspace{1em}

% ═══════════════════════════════════════════════════════════════
% ZEITPLANUNG
% ═══════════════════════════════════════════════════════════════

\abschnitt{Zeitplanung Tafelarbeit}

\begin{tabular}{|C{1.5cm}|L{12cm}|}
\hline
\textbf{Min} & \textbf{Aktivität} \\
\hline
0--5 & Wiederholung Kreuzregel (Beispiel MgO an Tafel) \\
\hline
5--15 & Aufgabe A gemeinsam: \#1--3 an Tafel, \#4--7 selbstständig \\
\hline
15--25 & Aufgabe B gemeinsam: \#2 (Na$_2$O) und \#3 (Cr$_2$O$_3$) an Tafel \\
\hline
25--30 & \#7 (Silberoxid) gemeinsam $\rightarrow$ Überleitung Reduktion \\
\hline
\end{tabular}

\end{document}
